\chapter{Introduzione}
\label{cap:introduzione}

Negli ultimi anni i sistemi IoT si sono diffusi anche nell'ambito della protezione
civile e della gestione delle emergenze, dove la disponibilità di informazioni
tempestive e affidabili può incidere direttamente sull'efficacia dei soccorsi.

Il terremoto del 2016 che ha colpito l'Italia centrale, e in particolare l'area di
Camerino, ha reso evidente questa esigenza: nelle prime ore successive a un evento
sismico, per i soccorritori è fondamentale sapere dove sono le persone, soprattutto
quando potrebbero trovarsi in un luogo difficilmente accessibile o non essere in grado
di muoversi. Da qui la domanda alla base della tesi: con quale
sensore è possibile rilevare una persona che non si muove?

In questo capitolo vengono presentate le motivazioni e gli obiettivi che hanno guidato
il lavoro, insieme a una panoramica della struttura della tesi. In particolare, la
\cref{sec:motivazioni} descrive il contesto progettuale e i motivi che hanno portato a
questo studio, la \cref{sec:obiettivo} ne definisce l'obiettivo e, infine, la
\cref{sec:struttura} ne illustra l'organizzazione.

\section{Motivazioni}
\label{sec:motivazioni}

Il progetto di tesi trova le sue motivazioni nel progetto
DIPME\footnote{Nella documentazione scientifica del gruppo di ricerca il progetto è
indicato anche con il nome \textbf{SAFE}; nel seguito si userà DIPME per
riferirsi allo stesso ecosistema.}, che ha come obiettivo la mitigazione del rischio
sismico attraverso l'integrazione di sensoristica IoT negli \emph{arredi} degli
ambienti scolastici~\cite{slideUprise}.

L'idea del progetto si articola su tre livelli.

\begin{description}[leftmargin=1.4em,style=nextline]
  \item[Arredi salva-vita] Banchi e pareti progettate dalla Scuola di Architettura
    e Design di UNICAM per resistere al crollo e offrire un volume di rifugio: il
    banco ha telaio doppio spaziale, piano rinforzato con lamiera forata
    antisfondamento, elemento di base dissipativo e connessione strutturale fra
    banchi adiacenti~\cite{pietroni2021furniture}.
  \item[Nodi di sensing] Nell'arredo è integrata l'elettronica della piattaforma
    \textbf{DIPME}, sviluppata con il partner industriale AM Microsystems. Il nodo
    dispiegato nel dimostratore scolastico~\cite{callisto2026dipme} è un'unità a
    batteria su microcontrollore STM32WLE con radio LoRa, un sensore \emph{PIR}
    (Panasonic EKMB139) per la presenza e sensori ambientali (BME280 per
    temperatura, umidità e pressione, SCD40 per la CO\textsubscript{2}); la scheda
    descritta nelle slide di progetto~\cite{slideUprise} prevede anche un sensore
    \emph{UWB}, un accelerometro e il Bluetooth. I nodi comunicano con un
    coordinatore LoRa (\textbf{DIPME-COORDINATOR}) e da questo con un gateway di
    bordo (\emph{edge}) su Linux.
  \item[Raccolta e visualizzazione] In emergenza la rete di terra non è
    disponibile: i messaggi LoRa (bande sub-GHz libere, 433/868~MHz in Europa,
    con portate dell'ordine di 3--5~km in ambiente urbano) vengono raccolti da
    \emph{droni} che ricostruiscono una mappa di calore geolocalizzata, consultabile
    offline dal tablet dei soccorritori. L'intero sistema è modellato e validato
    a priori sulla piattaforma di Digital Twin del gruppo di
    ricerca~\cite{callisto2023dt,unicamDT}.
\end{description}

Il sistema opera in due modalità: un ``tempo di pace'' di monitoraggio ordinario
e un ``tempo di guerra'' di emergenza sismica, attivato da un trigger (accelerometro
sul gateway, blackout della rete elettrica). Nel dimostratore descritto
in~\cite{callisto2026dipme} --- 42 nodi integrati nei banchi salva-vita di due aule
di una scuola secondaria delle Marche, in esercizio da ottobre 2025 a marzo 2026 ---
la transizione è comandata dal gateway al segnale di un accelerometro strutturale, e
in emergenza \emph{il bit di presenza decide la cadenza di trasmissione del nodo}: un
messaggio al minuto se il nodo rileva qualcuno, uno ogni quindici minuti se non rileva
nessuno. La qualità del rilevamento di presenza stabilisce quindi quali banchi vengono
ascoltati con priorità.

\todo{inserire una figura di sistema: arredo $\to$ DIPME $\to$ coordinator $\to$
gateway $\to$ drone $\to$ tablet}

\paragraph{Il limite del sensore adottato.}
Il DIPME-DEVICE rileva la presenza con un sensore PIR (\emph{Passive InfraRed}).
Il PIR è il sensore di presenza più diffuso al mondo per tre ragioni difficili
da battere: consuma qualche decina di microampere, costa pochi euro e non richiede
alcuna elaborazione. Ha tuttavia un limite che non dipende dalla qualità del
componente ma dal suo principio di funzionamento: l'elemento piroelettrico risponde
alla \emph{variazione} del flusso infrarosso incidente, non al suo valore. Una
persona ferma produce un flusso costante, quindi un segnale differenziale nullo:
\emph{nessun PIR, a nessun prezzo, può rilevare una persona immobile}
(\cref{sec:pir-principio}).

Nello scenario DIPME questo limite cade esattamente nel punto peggiore. Una
persona rifugiata sotto un banco dopo un crollo è, tipicamente, \textbf{ferma}:
per paura, per scelta, o perché non può muoversi. È il caso cieco del sensore
che il nodo monta oggi. Gli stessi autori della piattaforma indicano fra gli sviluppi
futuri l'estensione del livello di sensing con radar a onde millimetriche a
complemento del PIR, da valutare nei compromessi fra accuratezza, consumo, privacy e
complessità~\cite{callisto2026dipme}: questa tesi è, in concreto, quella valutazione.

I radar \emph{mmWave} a onda continua modulata in frequenza (FMCW) a 24~GHz
promettono di coprire proprio questo buco: rilevano i micro-movimenti
(fino, in linea di principio, all'oscillazione toracica del respiro), misurano la
distanza del bersaglio e non producono immagini, il che li rende compatibili con i
vincoli di privacy di un ambiente scolastico. Da qui le domande di ricerca della
tesi:

\begin{enumerate}[leftmargin=2em,itemsep=2pt]
  \item \emph{quanto} vale, in numeri, il vantaggio del mmWave sul PIR nel caso
        della persona ferma?
  \item quali \emph{dati} produce realmente un modulo mmWave di fascia consumer, e
        con quale accuratezza?
  \item a che \emph{costo} energetico, e con quale architettura di nodo?
  \item è possibile andare oltre la presenza binaria, verso una misura di
        \emph{vitalità} utile al triage?
\end{enumerate}

\section{Obiettivo}
\label{sec:obiettivo}

La tesi si propone di caratterizzare e confrontare sperimentalmente le due tecnologie
di rilevamento di presenza candidate per il nodo DIPME --- il sensore PIR
attualmente impiegato e il radar mmWave a 24~GHz --- e di stabilire, con misure
ripetibili e non con argomenti di principio, se e di quanto il radar sia necessario
nello scenario della persona immobile sotto l'arredo.

Attorno a questo obiettivo centrale il lavoro si articola in sei punti, concordati con
il relatore e qui riportati nella formulazione originale perché costituiscono la
traccia di tutto il documento.

\begin{enumerate}[label=\textbf{O\arabic*.},leftmargin=2.6em,itemsep=4pt]
  \item \textbf{Studio dei sensori} --- chi li usa, come vengono usati, quali dati
        producono (PIR e mmWave). \emph{Capitolo \ref{cap:sensori}.}
  \item \textbf{Analisi dei dati forniti} --- capire nel dettaglio quali grandezze
        i sensori mettono a disposizione: il PIR solo movimento sì/no, il mmWave
        micro-movimenti, distanza e potenzialmente respiro.
        \emph{Capitolo \ref{cap:moduli}.}
  \item \textbf{Comparazione con testing numerico} --- confronto sperimentale
        PIR vs mmWave con metriche quantitative. \emph{Capitolo \ref{cap:confronto}.}
  \item \textbf{Consumo energetico} --- a titolo informativo, sulla base dei
        datasheet, non essendo disponibile strumentazione di misura.
        \emph{Capitolo \ref{cap:consumi}.}
  \item \textbf{User interface web} --- l'ESP32 pubblica i dati, un sito web li
        visualizza in tempo reale, la stessa applicazione li salva con statistiche
        ed esporta in CSV per l'analisi numerica. \emph{\cref{sec:webui}, nel
        Capitolo \ref{cap:metodologia}.}
  \item \textbf{Indice di vitalità} --- obiettivo avanzato: oltre alla presenza
        sì/no, un algoritmo che quantifichi quanto una persona si muove
        (per esempio 50 su 100) e la classifichi di conseguenza.
        \emph{Capitolo \ref{cap:vitalita}.}
\end{enumerate}

\paragraph{Fuori scope.} Non fanno parte della tesi: la radio LoRa e il suo
protocollo, la raccolta con droni, il Digital Twin, la progettazione strutturale
degli arredi. Questi elementi compaiono solo nella \cref{sec:motivazioni} come
inquadramento e, dove pertinente, come vincolo di progetto (per esempio l'assenza
di connettività in emergenza, che orienta le scelte del \cref{cap:webui}).

\paragraph{Contributi.}
\label{sec:contributi}
Il perseguimento degli obiettivi appena elencati ha prodotto i seguenti contributi
originali:

\begin{itemize}[leftmargin=1.4em,itemsep=3pt]
  \item una \textbf{piattaforma di misura riproducibile} su ESP32 che acquisisce
        radar e PIR in parallelo a 5~Hz esatti con accesso ai dati per-gate del
        radar, con formato CSV unico condiviso da tutta la catena di analisi
        (\cref{sec:piattaforma});
  \item la \textbf{caratterizzazione quantitativa} del divario PIR/mmWave sulla
        persona ferma e sulla persona in movimento sul posto, con deviazioni
        standard su cinque ripetizioni per scenario (\cref{cap:confronto});
  \item la \textbf{calibrazione della misura di distanza} del LD2410B su 25 prove
        fra 1 e 5~m, con individuazione di un errore di scala sistematico e del
        coefficiente che lo corregge (\cref{sec:distanza});
  \item la \textbf{misura del respiro} con hardware consumer, verificata contro un
        ritmo imposto, e l'individuazione documentata di quale dei venti canali di
        energia disponibili sia effettivamente utilizzabile a ciascuna distanza
        (\cref{sec:respiro});
  \item la \textbf{quantificazione dell'attenuazione degli ostacoli} e della
        portata residua su entrambi i radar, cartongesso e porta chiusa compresi, e la
        misura della \textbf{portata massima} dei tre sensori nella stessa geometria
        (\cref{sec:ostacoli,sec:portata-massima});
  \item la \textbf{caratterizzazione del secondo radar} entro la portata
        dell'esemplare, con la taratura necessaria a farlo funzionare e i suoi limiti
        misurati (\cref{sec:ld2420-risultati});
  \item un \textbf{indice di vitalità} a tre classi, specificato, tarato e validato su
        trial separati, portato a bordo dell'ESP32 e verificato identico al calcolo di
        riferimento (\cref{cap:vitalita});
  \item un'\textbf{interfaccia web servita dal nodo} in Access Point, senza rete
        esterna, con esportazione CSV verificata equivalente alla catena seriale
        (\cref{cap:webui}).
\end{itemize}

\section{Struttura della Tesi}
\label{sec:struttura}

\begin{itemize}[leftmargin=1.4em,itemsep=4pt]
  \item \textbf{Sensori di presenza: principi e stato dell'arte}. Il primo capitolo
    offre una panoramica dei sensori di presenza, spiegando il principio fisico del
    PIR e quello del radar FMCW a onde millimetriche, e inquadrando entrambi rispetto
    alle alternative (UWB, ultrasuoni, visione). Presenta inoltre lo stato dell'arte
    sull'uso del mmWave per il monitoraggio dei parametri vitali senza contatto (O1).

  \item \textbf{I moduli in esame e i dati che producono}. Il secondo capitolo
    descrive in dettaglio i tre sensori impiegati --- il PIR HC-SR501 e i due radar
    Hi-Link LD2410B e LD2420 --- a partire dalla documentazione ufficiale del
    costruttore, con le grandezze che ciascuno mette a disposizione e i limiti
    dichiarati. Presenta poi la piattaforma di acquisizione realizzata su ESP32 e il
    formato dati unico usato da tutta la catena di analisi (O2).

  \item \textbf{Metodologia di confronto}. Il terzo capitolo fissa il metodo con cui i
    sensori sono stati confrontati: le geometrie e gli scenari, le ripetizioni e la
    ground truth, le metriche e le convenzioni di lettura dei dati, il protocollo
    ridotto adottato per il secondo radar. Presenta inoltre l'interfaccia web servita
    direttamente dall'ESP32 in modalità Access Point, senza alcuna rete esterna,
    realizzata a supporto della campagna con visualizzazione in tempo reale,
    statistiche ed esportazione CSV nello stesso formato della catena seriale (O5).

  \item \textbf{Confronto sperimentale}. Il quarto capitolo è il capitolo centrale
    della tesi e raccoglie i risultati della campagna di misura: accuratezza della
    distanza, falsi positivi e negativi, latenze di rilevamento e rilascio, copertura
    angolare, penetrazione degli ostacoli, portata massima e caratterizzazione del
    secondo radar. Ogni risultato è riportato con la sua dispersione su ripetizioni
    indipendenti (O3).

  \item \textbf{Indice di vitalità}. Il quinto capitolo misura il respiro con verifica
    a ritmo imposto e poi definisce, tara e valida un indice che quantifica quanto un
    soggetto si muove, oltre alla semplice presenza binaria, descrivendone il porting
    a bordo del microcontrollore e la verifica contro il calcolo di riferimento (O6).

  \item \textbf{Consumo energetico}. Il sesto capitolo analizza i consumi dei tre
    sensori a partire dai datasheet, ne stima l'autonomia a batteria nello scenario
    del nodo DIPME e discute l'architettura ibrida PIR + mmWave che ne discende (O4).

  \item \textbf{Conclusioni e sviluppi futuri}. La tesi si conclude rispondendo in
    sintesi alla domanda di partenza, dichiarando i limiti del lavoro svolto e
    indicando come può essere migliorato.
\end{itemize}

Le appendici raccolgono i dettagli del protocollo seriale dei due radar e i listati
principali del software sviluppato.
